%----------------------------------------------------------------------------------------
% VRAAGSTUKKEN. Alle opgaven zijn eigen werk.
%
% Opbouw: Probleem, Start, leerstof, en op een nieuwe pagina de Opdrachten en Test jezelf
% (2de graad: geen Koppel). Labels van dit hoofdstuk beginnen met h3-.
% Verschillen per pakket staan alleen in de oefeningenbank (pakketten:): een opdracht met kranen (5);
% de opdracht met de cijfers van een getal niet in 1.
%----------------------------------------------------------------------------------------

\mychapter{3}{Vraagstukken}

\begin{probleem}
Een kajakker vaart op een rivier 12~km stroomafwaarts. Dat duurt 1,5 uur. Daarna vaart hij dezelfde 12~km terug, tegen de stroom in, en dat duurt 3 uur. Hij peddelt de hele tijd even hard. Hoe snel zou hij varen in stilstaand water? En hoe snel stroomt de rivier?

\figuren{\begin{tikzpicture}[font=\scriptsize]
  \fill[accent!10] (0,0) rectangle (10,1.5);
  \draw[accent!60] (0,0) -- (10,0) (0,1.5) -- (10,1.5);
  \foreach \x in {1,2.7,6.4,8.1}{\draw[->,accent!70,thick] (\x,.75) -- ++(.9,0);}
  \node[text=accent] at (5,.75) {stroom};
  \draw[->,inkt,thick] (.4,1.15) -- (9.6,1.15);
  \draw[<-,inkt,thick,dashed] (.4,.35) -- (9.6,.35);
  \node[above] at (5,1.5) {stroomafwaarts: 12 km in 1,5 uur};
  \node[below] at (5,0) {stroomopwaarts: 12 km in 3 uur};
\end{tikzpicture}}
\end{probleem}

\lkr{Leerplan \pakketnaam: LPD \alleenvoor{1,2,3}{23}\alleenvoor{4,5}{27} en LPD 2, met LPD \alleenvoor{1,2,3}{22}\alleenvoor{4,5}{26} bij de twee tarieven (dit hoofdstuk: vraagstukken oplossen met een stelsel van twee eerstegraadsvergelijkingen met twee onbekenden, de oplossing controleren en interpreteren in de context, ook grafisch bij twee tarieven). LPD 2 vraagt vraagstukken oplossen door te mathematiseren en te demathematiseren: daarom eindigt elk voorbeeld met een antwoordzin en een controle in de tekst. Voorkennis: hoofdstuk 1 en 2 (grafisch en algebraïsch oplossen), en vraagstukken met één onbekende uit Eerstegraadsvergelijkingen; dat laatste wordt niet verondersteld. Laat leerlingen het probleem eerst proberen: velen rekenen $12 : 1{,}5 = 8$ en $12 : 3 = 4$ uit, maar weten dan niet hoe ze de stroom eruit halen. In sectie 3 wordt het probleem opgelost.}

Je kunt nu een stelsel oplossen, grafisch en algebraïsch. In dit hoofdstuk gebruik je dat om vraagstukken op te lossen. Het moeilijkste deel is niet het rekenen, maar het opstellen: de tekst omzetten in twee vergelijkingen. Je leert hoe je twee onbekenden kiest, hoe je in de tekst twee voorwaarden zoekt, en hoe je het antwoord controleert. Daarna zie je de soorten vraagstukken die het vaakst voorkomen: prijzen, mengsels, procenten, snelheden, meetkunde en tarieven.

\begin{oefeningen}[Start]
\oef{h3-start-1}

\oef{h3-verken}
\end{oefeningen}

%========================================================================================
\section{Van vraagstuk naar stelsel}

\subsection{Twee onbekenden, twee voorwaarden}

In de Verken zocht je door te proberen hoeveel leerlingen en hoeveel volwassenen een ticket kochten. Proberen lukt, maar het duurt lang, en bij moeilijke getallen vind je zo het antwoord nooit. Er zijn twee getallen die je niet kent: het aantal leerlingen en het aantal volwassenen. Geef ze elk een letter. De tekst zegt twee dingen over die getallen: hoeveel tickets er samen zijn, en hoeveel ze samen opbrachten. Elke voorwaarde wordt een vergelijking, en samen vormen ze een stelsel.

\voorbeeld{}{Op een schoolfeest kost een ticket 6 euro voor een leerling en 9 euro voor een volwassene. Er werden 120 tickets verkocht, en de kassa telde 870 euro. Hoeveel leerlingen en hoeveel volwassenen kochten een ticket?\label{vb:h3-tickets}}{%
\uitleg{We kiezen voor elk onbekend aantal een letter, en zetten de twee voorwaarden uit de tekst om in vergelijkingen. Dan lossen we het stelsel op met substitutie, want in de eerste vergelijking staat $x$ zonder getal ervoor.}
Kies de onbekenden en zeg wat ze betekenen: $x$ is het aantal leerlingen en $y$ het aantal volwassenen.
\begin{stappen}
\stelsel{x + y &= 120}{6x + 9y &= 870} & Samen 120 tickets. De leerlingen betalen $6x$ euro, de volwassenen $9y$ euro. \\[3ex]
x = 120 - y & Los de eerste vergelijking op naar $x$. \\
6(120 - y) + 9y = 870 & Substitueer in de tweede vergelijking. \\
720 + 3y = 870 \ \Rightarrow \ y = 50 & Werk uit: $-6y + 9y = 3y$, en $3y = 150$. \\
x = 120 - 50 = 70 & Vul $y = 50$ in.
\end{stappen}
Controle in de tekst: $70 + 50 = 120$ tickets, en $70 \cdot 6 + 50 \cdot 9 = 420 + 450 = 870$ euro. \textit{Er kwamen 70 leerlingen en 50 volwassenen.}}

Controleer altijd in de tekst, niet in je stelsel. Heb je een vergelijking verkeerd opgesteld, dan klopt je oplossing wel in het stelsel, maar niet in het vraagstuk. Kijk ook of het antwoord kan: $70{,}5$ leerlingen of een negatieve prijs is geen antwoord, en wijst op een fout.

\begin{werkwijze}
Lees de tekst en zeg wat je zoekt. Kies twee onbekenden $x$ en $y$, en schrijf erbij wat ze betekenen, met de eenheid. Zoek in de tekst twee voorwaarden en zet elke voorwaarde om in een vergelijking. Los het stelsel op met de methode die het best past. Controleer in de tekst, en antwoord met een zin.
\end{werkwijze}

\nujij{Een tablet en een hoes kosten samen 380 euro. De tablet kost 300 euro meer dan de hoes. Hoeveel kost de hoes?}{$t + h = 380$ en $t = h + 300$ geven $2h + 300 = 380$, dus de hoes kost 40 euro. Wie 80 euro antwoordt, vergeet dat de tablet dan $380$ euro kost, en samen $460$.}

\begin{oefeningen}
\oef{h3-opstellen}
\end{oefeningen}

\subsection{Eén onbekende of twee?}

Kon het vraagstuk van het schoolfeest ook met één onbekende? Ja. Zijn er $x$ leerlingen, dan zijn er $120 - x$ volwassenen, en de opbrengst geeft de vergelijking $6x + 9(120 - x) = 870$. Zo los je in het handboek Eerstegraadsvergelijkingen een vraagstuk op met twee soorten, zoals witte en grijze broden. Kijk goed: die vergelijking is precies wat je krijgt na de substitutie in voorbeeld~\ref{vb:h3-tickets}. Met één onbekende doe je de substitutie dus in je hoofd, nog voor je iets opschrijft.

Bij eenvoudige vraagstukken maakt dat weinig uit. Maar hoe meer voorwaarden er in de tekst staan, hoe moeilijker het wordt om alles in één letter uit te drukken. Met twee onbekenden schrijf je elke zin apart om, en het rekenwerk komt pas daarna.

\voorbeeld{}{Een getal is 7 minder dan het dubbel van een ander getal. Samen zijn de twee getallen 41. Welke getallen zijn het?\label{vb:h3-getallen}}{%
\uitleg{We zetten elke zin apart om in een vergelijking. De eerste vergelijking zegt al waaraan $x$ gelijk is, dus substitutie ligt voor de hand.}
Noem het eerste getal $x$ en het tweede getal $y$.
\begin{stappen}
\stelsel{x &= 2y - 7}{x + y &= 41} & Het dubbel van het tweede getal, verminderd met 7. Samen 41. \\[3ex]
(2y - 7) + y = 41 & Substitueer $x$ in de tweede vergelijking. \\
3y = 48 \ \Rightarrow \ y = 16 & Breng $-7$ over en deel door 3. \\
x = 2 \cdot 16 - 7 = 25 & Vul in.
\end{stappen}
Controle in de tekst: het dubbel van 16 is 32, en $32 - 7 = 25$. Samen: $25 + 16 = 41$. \textit{De getallen zijn 25 en 16.} Met één onbekende noem je het tweede getal $y$ en het eerste meteen $2y - 7$: dat is de tweede regel van de berekening.}

\begin{oefeningen}
\oef{h3-getallen}
\end{oefeningen}

%========================================================================================
\section{Prijzen, mengsels en procenten}

Veel vraagstukken gaan over twee soorten van iets: twee producten met elk een eigen prijs, twee soorten koffie in één mengsel, twee sappen met een ander percentage fruit. Hoe zet je zo'n tekst om? Het idee is telkens hetzelfde: een aantal maal een prijs geeft een bedrag, en bedragen mag je optellen.

\subsection{Prijzen en aantallen}

Soms ken je de prijzen niet, maar wel wat twee verschillende bestellingen samen kosten. Elke bestelling geeft dan een vergelijking.

\voorbeeld{}{In een bakkerij kosten 3 croissants en 2 sandwiches samen 9,30 euro. Voor 5 croissants en 1 sandwich betaal je 8,50 euro. Wat kost een croissant, en wat kost een sandwich?\label{vb:h3-prijzen}}{%
\uitleg{De twee prijzen zijn de onbekenden, en elke bestelling geeft een vergelijking. In de tweede vergelijking staat $y$ zonder getal ervoor: die lossen we op naar $y$.}
Stel: een croissant kost $x$ euro en een sandwich $y$ euro.
\begin{stappen}
\stelsel{3x + 2y &= 9{,}30}{5x + y &= 8{,}50} & De twee bestellingen. \\[3ex]
y = 8{,}50 - 5x & Los de tweede vergelijking op naar $y$. \\
3x + 2(8{,}50 - 5x) = 9{,}30 & Substitueer in de eerste. \\
17 - 7x = 9{,}30 \ \Rightarrow \ x = 1{,}10 & Werk uit: $7x = 7{,}70$. \\
y = 8{,}50 - 5{,}50 = 3 & Vul in.
\end{stappen}
Controle: $3 \cdot 1{,}10 + 2 \cdot 3 = 3{,}30 + 6 = 9{,}30$ en $5 \cdot 1{,}10 + 3 = 8{,}50$. \textit{Een croissant kost 1,10 euro, een sandwich 3 euro.}}

\begin{oefeningen}
\oef{h3-prijzen}
\end{oefeningen}

\subsection{Mengsels}

Een koffiebrander mengt twee soorten bonen. Wat kost een kilogram van het mengsel? Dat hangt af van hoeveel hij van elke soort neemt. Bij een mengsel zijn er altijd twee voorwaarden: de hoeveelheden samen geven de hoeveelheid van het mengsel, en de waarde van de delen samen geeft de waarde van het mengsel. Door te mengen gaat er geen geld verloren en komt er geen bij.

\voorbeeld{}{Een koffiebrander mengt bonen van 24 euro per kg met bonen van 15 euro per kg. Hij wil 30~kg mengsel dat 21 euro per kg kost. Hoeveel kg van elke soort heeft hij nodig?\label{vb:h3-mengsel}}{%
\uitleg{De onbekenden zijn de twee hoeveelheden. De waarde van een hoeveelheid is het aantal kg maal de prijs per kg. De waarde van het hele mengsel is $30 \cdot 21 = 630$ euro.}
Stel: hij neemt $x$~kg bonen van 24 euro en $y$~kg bonen van 15 euro.
\begin{stappen}
\stelsel{x + y &= 30}{24x + 15y &= 630} & Samen 30~kg. De waarde van de delen is samen de waarde van het mengsel. \\[3ex]
24(30 - y) + 15y = 630 & Substitueer $x = 30 - y$. \\
720 - 9y = 630 \ \Rightarrow \ y = 10 & Werk uit: $9y = 90$. \\
x = 30 - 10 = 20 & Vul in.
\end{stappen}
Controle: $20 \cdot 24 + 10 \cdot 15 = 480 + 150 = 630$ euro voor 30~kg, dus 21 euro per kg. \textit{Hij mengt 20~kg van de dure bonen met 10~kg van de goedkope.}}

Kan dit kloppen? De prijs van het mengsel, 21 euro, ligt dichter bij 24 euro dan bij 15 euro. Er moet dus meer van de dure soort in. Zo'n schatting vooraf toont snel of je antwoord kan.

\begin{oefeningen}
\oef{h3-mengsel}
\end{oefeningen}

\subsection{Procenten}

Ook bij een percentage gaat het om een deel van een geheel. Een siroop met 10\,\% suiker bevat in 1 liter $0{,}1$ liter suiker, in $x$ liter dus $0{,}1x$ liter. De tweede voorwaarde gaat dan over de hoeveelheid suiker: die verandert niet als je mengt.

\voorbeeld{}{Een ijsmaker heeft siroop met 10\,\% suiker en siroop met 30\,\% suiker. Ze wil 2 liter siroop met 25\,\% suiker. Hoeveel liter neemt ze van elke soort?\label{vb:h3-procent}}{%
\uitleg{De onbekenden zijn de twee hoeveelheden siroop. De suiker in de twee delen is samen de suiker in het mengsel: 25\,\% van 2 liter is $0{,}5$ liter.}
Stel: ze neemt $x$ liter siroop met 10\,\% suiker en $y$ liter met 30\,\%.
\begin{stappen}
\stelsel{x + y &= 2}{0{,}1x + 0{,}3y &= 0{,}5} & Samen 2 liter. De suiker uit de twee delen is samen $0{,}5$ liter. \\[3ex]
\stelsel{x + y &= 2}{x + 3y &= 5} & Vermenigvuldig de tweede vergelijking met 10: dan zijn de kommagetallen weg. \\[3ex]
2y = 3 \ \Rightarrow \ y = 1{,}5 & Combinatie: trek de eerste vergelijking af van de tweede. \\
x = 2 - 1{,}5 = 0{,}5 & Vul in.
\end{stappen}
Controle: $0{,}1 \cdot 0{,}5 + 0{,}3 \cdot 1{,}5 = 0{,}05 + 0{,}45 = 0{,}5$ liter suiker in 2 liter, dat is 25\,\%. \textit{Ze neemt $0{,}5$ liter van de siroop met 10\,\% en $1{,}5$ liter van die met 30\,\%.}}

Een veelgemaakte fout is de percentages zelf optellen: $10\,\% + 30\,\%$ is geen 25\,\%. Percentages van verschillende hoeveelheden mag je niet zomaar optellen, de hoeveelheden suiker wel.

\begin{oefeningen}
\oef{h3-procent}
\end{oefeningen}

%========================================================================================
\section{Beweging}

Hoe ver kom je, als je weet hoe snel en hoe lang je onderweg bent? Uit de lessen fysica weet je dat afstand gelijk is aan snelheid maal tijd. Rij je 2 uur aan 50~km/h, dan leg je $50 \cdot 2 = 100$~km af. Omgekeerd vind je de snelheid door de afstand te delen door de tijd. Met die ene regel los je alle vraagstukken in deze sectie op.

\subsection{Met en tegen de stroom}

Nu kunnen we het probleem van de kajakker oplossen. Op een rivier vaart een boot niet alleen met zijn eigen snelheid: het water neemt hem ook mee. Stroomafwaarts tellen de twee snelheden op, stroomopwaarts werkt de stroom tegen en trek je ze af.

\voorbeeld{}{Terug naar het probleem. Een kajakker vaart 12~km stroomafwaarts in 1,5 uur, en dezelfde 12~km terug in 3 uur. Hoe snel vaart hij in stilstaand water, en hoe snel stroomt de rivier?\label{vb:h3-terug}}{%
\uitleg{We zoeken twee snelheden: die van de kajak in stilstaand water en die van de stroom. Voor elke rit zetten we de afstand, de snelheid en de tijd in een tabel.}
Stel: $v$ is de snelheid van de kajak in stilstaand water en $w$ de snelheid van de stroom, in km/h.
\naast{\begin{stappenvrij}{3.4cm}{5cm}
\stelsel{1{,}5\,(v + w) &= 12}{3\,(v - w) &= 12} & Snelheid maal tijd is de afstand, voor elke rit. \\[3ex]
\stelsel{v + w &= 8}{v - w &= 4} & Deel de eerste vergelijking door 1,5 en de tweede door 3. \\[3ex]
2v = 12 \ \Rightarrow \ v = 6 & Combinatie: tel de vergelijkingen op. \\
w = 8 - 6 = 2 & Vul in.
\end{stappenvrij}}{{\renewcommand{\arraystretch}{1.4}\small
\begin{tabular}{|c|c|c|c|}
\hline
 & \textbf{afstand} & \textbf{snelheid} & \textbf{tijd} \\ \hline
mee & $12$ & $v + w$ & $1{,}5$ \\ \hline
tegen & $12$ & $v - w$ & $3$ \\ \hline
\end{tabular}}}
Controle: stroomafwaarts vaart hij $6 + 2 = 8$~km/h, en $8 \cdot 1{,}5 = 12$~km. Stroomopwaarts vaart hij $6 - 2 = 4$~km/h, en $4 \cdot 3 = 12$~km. \textit{De kajakker vaart 6~km/h in stilstaand water, en de rivier stroomt 2~km/h.}}

Het gemiddelde van 8 en 4 is dus de snelheid in stilstaand water, en de stroom is de helft van het verschil. Dat is geen toeval: tel je $v + w$ en $v - w$ op, dan valt $w$ weg. Hetzelfde geldt voor een vliegtuig met de wind mee of tegen.

\begin{werkwijze}
Maak een tabel met voor elke beweging de afstand, de snelheid en de tijd, in dezelfde eenheden. Gebruik afstand $=$ snelheid $\cdot$ tijd om een vergelijking per beweging op te stellen. Met de stroom of de wind mee tel je de snelheden op, ertegen trek je ze af.
\end{werkwijze}

\subsection{Naar elkaar toe}

Wat als twee mensen tegelijk vertrekken en naar elkaar toe rijden? Elk uur komen ze dan dichter bij elkaar, met de som van hun snelheden. Op het moment dat ze elkaar ontmoeten, hebben ze samen de hele afstand afgelegd.

\voorbeeld{}{Twee fietsers wonen 60~km van elkaar. Ze vertrekken tegelijk en fietsen naar elkaar toe. Na 1,5 uur ontmoeten ze elkaar. De ene fietst 4~km/h sneller dan de andere. Hoe snel fietst elk?\label{vb:h3-tegemoet}}{%
\uitleg{De onbekenden zijn de twee snelheden. Samen leggen ze in 1,5 uur 60~km af.}
Stel: de snelste fietst $x$~km/h en de traagste $y$~km/h.
\begin{stappen}
\stelsel{1{,}5x + 1{,}5y &= 60}{x - y &= 4} & De afstanden samen zijn 60~km. Het verschil in snelheid is 4~km/h. \\[3ex]
\stelsel{x + y &= 40}{x - y &= 4} & Deel de eerste vergelijking door 1,5. \\[3ex]
2x = 44 \ \Rightarrow \ x = 22, \ y = 18 & Tel op, en vul in.
\end{stappen}
Controle: in 1,5 uur fietst de ene $22 \cdot 1{,}5 = 33$~km, de andere $18 \cdot 1{,}5 = 27$~km, samen 60~km. \textit{De ene fietst 22~km/h, de andere 18~km/h.}}

\nujij{Een vliegtuig vliegt 1500~km met de wind mee in 2 uur, en terug in $2{,}5$ uur. Hoe snel vliegt het bij windstil weer?}{$v + w = 750$ en $v - w = 600$, dus $v = 675$~km/h (en de wind waait 75~km/h).}

\begin{oefeningen}
\oef{h3-beweging}
\end{oefeningen}

%========================================================================================
\section{Meetkunde en tarieven}

\subsection{Omtrek en hoeken}

Ook in de meetkunde zoek je vaak twee lengtes of twee hoeken tegelijk. Je hebt dan formules nodig die je al kent: de omtrek van een rechthoek is twee keer de som van lengte en breedte, en de hoeken van een driehoek zijn samen $180^\circ$. Een schets met de letters erbij helpt om de voorwaarden te zien.

\voorbeeld{}{Een rechthoek heeft een omtrek van 38~cm. Maak je de lengte 3~cm korter en de breedte 2~cm langer, dan krijg je een vierkant. Hoe lang en hoe breed is de rechthoek?\label{vb:h3-rechthoek}}{%
\uitleg{De onbekenden zijn de lengte en de breedte. Een vierkant heeft vier gelijke zijden: na de verandering zijn lengte en breedte dus gelijk.}
Stel: de rechthoek is $x$~cm lang en $y$~cm breed.
\naast{\begin{stappenvrij}{5.2cm}{4cm}
\stelsel{2(x + y) &= 38}{x - 3 &= y + 2} & De omtrek. De nieuwe lengte is gelijk aan de nieuwe breedte. \\[3ex]
\stelsel{x + y &= 19}{x - y &= 5} & Deel door 2, en breng over. \\[3ex]
2x = 24 \ \Rightarrow \ x = 12, \ y = 7 & Tel op, en vul in.
\end{stappenvrij}}{\begin{tikzpicture}[scale=.2,font=\scriptsize]
  \draw[accent,thick,fill=accent!10] (0,0) rectangle (12,7);
  \draw[inkt,dashed] (0,0) rectangle (9,9);
  \node[below] at (6,0) {$x$};
  \node[left] at (0,3.5) {$y$};
  \node[above,text=inkt] at (4.5,9) {vierkant};
\end{tikzpicture}}
Controle: $2 \cdot (12 + 7) = 38$~cm, en $12 - 3 = 9 = 7 + 2$. \textit{De rechthoek is 12~cm lang en 7~cm breed.}}

\voorbeeld{}{In een rechthoekige driehoek is de ene scherpe hoek $18^\circ$ groter dan de andere. Hoe groot zijn de twee scherpe hoeken?\label{vb:h3-hoeken}}{%
\uitleg{De rechte hoek is $90^\circ$, dus de twee scherpe hoeken zijn samen $180^\circ - 90^\circ = 90^\circ$.}
Noem de kleinste scherpe hoek $x$ en de grootste $y$, in graden.
\begin{stappen}
\stelsel{x + y &= 90}{y &= x + 18} & Samen $90^\circ$. De ene is $18^\circ$ groter. \\[3ex]
x + (x + 18) = 90 \ \Rightarrow \ x = 36 & Substitueer: $2x = 72$. \\
y = 36 + 18 = 54 & Vul in.
\end{stappen}
Controle: $36^\circ + 54^\circ + 90^\circ = 180^\circ$. \textit{De scherpe hoeken zijn $36^\circ$ en $54^\circ$.}}

\begin{oefeningen}
\oef{h3-meetkunde}
\end{oefeningen}

\subsection{Twee tarieven vergelijken}

Welk abonnement is het voordeligst? Dat hangt meestal af van hoe vaak je het gebruikt. Bij elk tarief hoort een eerstegraadsfunctie: de prijs $y$ in functie van het gebruik $x$, met een vaste prijs $q$ en een prijs $m$ per keer. Elk voorschrift $y = mx + q$ is een vergelijking met twee onbekenden, en samen vormen ze een stelsel. De oplossing is het snijpunt van de twee rechten: daar betaal je bij beide evenveel. Links en rechts van het snijpunt zie je op de grafiek welk tarief goedkoper is.

\voorbeeld{}{Een zwembad heeft twee formules. Met de beurtenkaart betaal je 4 euro per beurt. Met het jaarabonnement betaal je 60 euro per jaar, en daarbovenop 1,50 euro per beurt. Vanaf hoeveel beurten per jaar is het abonnement voordeliger?\label{vb:h3-tarief}}{%
\uitleg{We schrijven de jaarprijs bij elke formule als een vergelijking, tekenen beide rechten en berekenen het snijpunt exact.}
Stel: $x$ is het aantal beurten per jaar en $y$ de jaarprijs in euro.
\naast{\begin{stappenvrij}{4.4cm}{3.5cm}
\stelsel{y &= 4x}{y &= 60 + 1{,}5x} & De beurtenkaart en het abonnement. \\[3ex]
4x = 60 + 1{,}5x & Substitueer $y = 4x$. \\
2{,}5x = 60 \ \Rightarrow \ x = 24 & Breng over en deel door 2,5. \\
y = 4 \cdot 24 = 96 & Het snijpunt is $(24, 96)$.
\end{stappenvrij}}{\begin{tikzpicture}[x=.13cm,y=.016cm,font=\scriptsize]
  \roosterleeg[$x$][$y$]{0}{5}{40}{0}{20}{160}
  \xgetallen{10,20,30,40}\ygetallen{40,80,120,160}\oorsprong
  \rechte{0}{40}{4}{0}
  \rechte[tweede]{0}{40}{1.5}{60}
  \node[accent,left] at (30,140) {kaart};
  \node[inkt,below] at (33,100) {abonnement};
  \punt{24}{96}\aflezen{24}{96}
\end{tikzpicture}}
Bij 24 beurten betaal je bij beide 96 euro. Op de grafiek ligt de rechte van het abonnement rechts van het snijpunt onder die van de beurtenkaart. \textit{Zwem je meer dan 24 keer per jaar, dan is het abonnement voordeliger; minder dan 24 keer, dan de beurtenkaart.} Controle met 30 beurten: de kaart kost 120 euro, het abonnement $60 + 45 = 105$ euro.}

Het snijpunt alleen is dus nog geen antwoord. De vraag is meestal \textit{wanneer} een tarief voordeliger is, en dat lees je af op de grafiek, links en rechts van het snijpunt. Controleer met een getal aan één kant van het snijpunt.

\begin{oefeningen}
\oef{h3-tarief}
\end{oefeningen}

%========================================================================================
\opdrachtenkop

\begin{oefeningen}
\oef{h3-o1}

\oef{h3-o2}

\oef{h3-o3}

\oef{h3-o4}

\oef{h3-o5}

\oef{h3-o6}

\oef{h3-o7}

\oef{h3-o8}

\oef{h3-o9}

\oef{h3-o10}

\oef{h3-o11}

\oef{h3-o12}

\oef{h3-o13}

\oef{h3-o14}

\oef{h3-o15}

\oef{h3-o16}

\oef{h3-o17}

\oef{h3-o18}

\oef{h3-o19}

\oef{h3-o20}

\oef{h3-o21}

\oef{h3-o22}


\oef{h3-o24}

\oef{h3-o25}

\oef{h3-o26}
\end{oefeningen}

%----------------------------------------------------------------------------------------
\needspace{14\baselineskip}
\opdrachtsectie{Test jezelf}

\begin{oefeningen}
\oef{h3-t1}

\oef{h3-t2}

\oef{h3-t3}

\oef{h3-t4}

\oef{h3-t5}

\oef{h3-t6}

\oef{h3-t7}
\end{oefeningen}

\herhaalschema{%
  \fv{oef:h3-t1}{vb:h3-getallen}{oef:h3-getallen}%
  \fv{oef:h3-t2}{vb:h3-prijzen}{oef:h3-prijzen}%
  \fv{oef:h3-t3}{vb:h3-mengsel}{oef:h3-mengsel}%
  \fv{oef:h3-t4}{vb:h3-terug}{oef:h3-beweging}%
  \fv{oef:h3-t5}{vb:h3-rechthoek}{oef:h3-meetkunde}%
  \fv{oef:h3-t6}{vb:h3-tarief}{oef:h3-tarief}%
  \fv{oef:h3-t7}{vb:h3-procent}{oef:h3-procent}}

\ifzelfstudie
\section*{Maak je eigen samenvatting}
Sluit dit boek en schrijf in je schrift: de stappen om een vraagstuk met een stelsel op te lossen, waarom je in de tekst controleert, en voor elke soort vraagstuk de twee voorwaarden die je meestal vindt: prijzen en aantallen, mengsels, procenten, met en tegen de stroom, naar elkaar toe, omtrek en hoeken, en twee tarieven. Zeg bij twee tarieven ook wat het snijpunt betekent. Vergelijk daarna met de voorbeelden.
\fi
